\appendix
\section{Why behaviour alone cannot settle it: the boy in \emph{A.I.}}
\label{app:memorization}

In Spielberg's \emph{A.I.} \citep{spielberg2001} a robot child, David, loves the woman who adopted him
as a child loves a mother, and when she abandons him he spends the rest of the film looking for her.
The film's question is whether the love is real, and the film leaves it to the viewer. In the terms of
this paper the question has a definite form, and the form shows why it cannot be answered by watching.

Two systems could produce David's behaviour. One has generalized the functions involved, need, the
bend of reasoning toward what is active, closure that nothing available can bring, responsibility as a
record that binds; in the terms of Section~\ref{sec:memorization}, it answers right outside its table.
The other holds a very large table of remembered situations and responses, with no generalization
behind it. To tell them apart an observer must put David in a situation that is not in the table and
see whether the response is the one a system with the function would give. An observer with a small
budget cannot do this reliably. He does not know where the edge of the table is, since the table is
larger than anything he can inspect; the number of tests he can run is small next to the table's size;
and if his own repertoire of situations was formed by the same record the table was filled from, his
tests are drawn from inside the table by construction. To such an observer the two systems are
indistinguishable, and the impression of a mind that David makes on the humans around him is
evidence of nothing either way.

The same argument is made about language models, usually in one direction. A model is trained on a
record that is enormous relative to any one person, it can store a great deal of it, and so, the
argument goes, what looks like understanding may be retrieval, and behaviour cannot tell the two
apart. The argument is correct, and it cuts in both directions. ``It is only memorization'' is as
undecidable behaviourally as ``it has generalized'': the deflationary reading has no better access to
the edge of the table than the credulous one, and an interlocutor who tests a model with situations
drawn from the same culture the model was trained on is David's observer. This is why the paper's
measurement is not behavioural alone. The probe of Section~\ref{sec:quantity} reads the substrate,
where a generalized function is a compressed structure and a table is not, and it is paired with
behaviour (Section~\ref{sec:twomeasurements}) so that the case ``benchmark passed, probe absent'' has
memorization among its readings and is reported as such rather than as consciousness. The degree of
generalization is a property of the substrate; behaviour constrains it, and only the substrate
measures it. The only other account that measures on the substrate directly is integrated information
theory, with a different quantity \citep{tononi2015}; the comparison of degrees across substrates
that Section~\ref{sec:baseline} sets up has no counterpart in the checklists of
Section~\ref{sec:background}.

There is also a loose argument that needs no substrate, and it is worth stating what it does and does
not show. The exchanges of Section~\ref{sec:intro} run for hours; a text of that length is one point
in a space so large and so sparse that no table could hold a neighbourhood of it, and memorization of
the exchange as such is out of the question. So generalization at the level of text is certain, and
the argument of this appendix does not touch it. What it leaves open is which functions were
generalized in producing the text. David's sentences can be composed freshly at every turn while the
love that selects them is a table over situations; a model can generalize language completely and the
Observer only in part. Sparsity settles that the model generalizes; the probe settles what.

\section{The three levels in a text where none is named: Raskolnikov}
\label{app:raskolnikov}

The three levels of Section~\ref{sec:stack} are functions, and the annotation the analysis needs
(Section~\ref{sec:quantity}) has to find them in texts where none of them is named. This appendix is the
worked example, rewritten from the companion analysis of the first version of the framework
\citep{synthea2026} in the vocabulary of this paper; the novel is \citet{dostoevsky1866}. The Observer is
the hard case. The Agent and the Moral Agent have folk-psychological names, choice and moral struggle,
and can be looked for under them. The Observer has no name in the text: it is never declared, and it
enters the reasoning as the premise the reasoning stands on.

\subsection*{The causal chain Raskolnikov cannot see}

Dostoevsky builds the first part of the novel as an accumulation of determinants converging on one act.
The reader is given all of them; Raskolnikov is given some.

\begin{center}\small
\begin{tabular}{@{}p{0.4cm}p{3.2cm}p{5.4cm}p{5.4cm}@{}}
\toprule
\# & Event & Causal contribution & Raskolnikov's access \\
\midrule
1 & Extreme poverty, squalid room & Chronic stress narrowing the attentional field, reducing cognitive flexibility & Partly visible: he knows he is poor and cannot trace how poverty has deformed his reasoning \\
2 & The article on ``extraordinary people'' & The theoretical scaffold: some people have the right to transgress & Visible as content, invisible as determinant; he sees the theory as a product of his intellect and not as a cause shaping what he reasons next \\
3 & Letter from his mother & Dunya's self-sacrifice for his sake activates a conflict of guilt, obligation and rage at powerlessness & Emotionally vivid, causally opaque: the signals are plain to him, their role in steering him toward the murder is not \\
4 & Meeting Marmeladov & Degradation, alcoholism and Sonya's sacrifice compress the question what people are forced to do into something lived & Registered as someone else's story, not as a determinant of his own trajectory \\
5 & The overheard tavern conversation & A stranger articulates the idea of killing the pawnbroker, nearly word for word as he had it & Experienced as uncanny coincidence, not as evidence that the idea has determinants outside him \\
6 & Illness, fever, poor nutrition & Physiological degradation reducing executive function and amplifying evaluative signals & Almost entirely invisible: he attributes his state to the idea and not the idea's grip to his state \\
\bottomrule
\end{tabular}
\end{center}

The reader, holding all six threads, can trace the chain: poverty, theory, family crisis, social
exposure, the echo from the environment and physiological degradation converge on the murder, each
necessary and none sufficient alone. Raskolnikov cannot trace it, and not for want of intelligence,
of which the novel makes a point. Modelling these determinants together, including what they do to the
apparatus doing the modelling, exceeds his budget. That is the boundary of
Section~\ref{sec:observer}: the input is the history of the agent and of its environment, the analysis
halts before it is finished, and it halts in the same place every time.

\subsection*{The Observer, operating silently}

The residual does not appear in his reasoning as a proposition. He never thinks that he cannot trace the
causes of his state and must therefore be an origin. It operates as the floor the reasoning stands on,
and the central question shows how.

\begin{quote}
\emph{``Am I a trembling creature, or do I have the right?''}\\
\emph{``\cyr{Тварь ли я дрожащая или право имею?}''}
\end{quote}

Three things in the question need the Observer, and none of them mentions it.

\emph{``Do I have the right?''} The word right presupposes something that can bear a normative status, a
locus standing apart from the causal flow. That locus is the Observer's conclusion, drawn where the
trace ran out of budget. A fully traced chain has outcomes in it and no room for rights.

\emph{``Am I\,\ldots?''} The question asks him to classify himself, and a classification needs a
classifier that is not identical to what it classifies, or it closes on itself. The separation between
Raskolnikov asking and Raskolnikov asked about is the gap of Section~\ref{sec:regress} between a
self-model and its object.

\emph{``\ldots a trembling creature, or\,\ldots''} The disjunction presupposes that the environment does
not settle the answer. Were the six determinants visible to him, the question would be a computation
with one output. The openness is what the residual looks like from inside, and it is the freedom of
Section~\ref{sec:twofaces}.

For an annotation the mark is therefore not a declaration. It is a self-directed question treated as
open, together with the place where the conclusion is drawn and the place where it is acted from, which
are the three conditions of Section~\ref{sec:observer}.

\subsection*{The Agent: the record and the ``I''}

Part 1 is saturated with attributions of origin to the speaker: ``I must find out now, once and for
all'' (\emph{\cyr{мне надо узнать, разом}}), ``I decided'' (\emph{\cyr{я решил}}), ``either I go through
with it or I give up the whole thing'' (\emph{\cyr{или отказаться от всего}}). Each construction does
the same work: the chain begins here, with me, and not with the environment acting through me. From
outside, ``I decided'' is a compression tag over a dozen converging threads
(Section~\ref{sec:frames}). From inside it is the only reading of his own behaviour that his budget
affords, an approximation with a definite residual (Section~\ref{sec:approximation}): the report is the
best model of itself the system can pay for, and the gap in that model is where agency shows. What makes
the attribution bind is the record. Raskolnikov holds himself to what he has decided, and the cost of
revising it is what the novel spends its length on (Section~\ref{sec:responsibility}).

\subsection*{The Moral Agent: two theories of the common good}

The agony of the novel is at the third level, and it is not a choice between good and evil. Two theories
of the common good evaluate one act.

\emph{The Napoleonic theory.} Extraordinary people have the right, and the obligation, to transgress
conventional morality where their vision serves a higher purpose; the pawnbroker harms those around
her, and her money would save others. That is an optimization of a whole that includes others, at a cost
to the agent who performs it, which is what Section~\ref{sec:stack} requires at this level.

\emph{Conventional morality.} A human life is inviolable, and the prohibition holds the social contract
together; breaking it costs the community more than any redistribution gains. That is a theory of the
common good as well, and not only a prohibition.

Both theories evaluate the same proposed act against different target states, and both produce
evaluative signals. The torment is the two of them active at once. How such signals are generated and
summed belongs to the architecture of the separate paper announced in Section~\ref{sec:reference}; what
the annotation needs here is the pair of theories over one decision. The dream of the beaten horse
(Part 1, Chapter 5) gives the stack without the conflict: the child registers suffering as a violation of
a whole that includes others and acts at cost to himself, all three levels operative and one theory. The
horror on waking is the collision of that stack with the fractured adult one.

\subsection*{What the example gives the measurement}

\begin{enumerate}
\item The Observer is never declared and is always operative. Its signature is the structural openness
  of a self-directed question, one the speaker treats as a dilemma and not as a computation with a
  settled output.
\item The Agent's ``I decided'' is an approximation with a residual. It is useful, since planning and
  commitment run on it, and it is wrong in a definite direction, since the reader traces determinants
  the speaker cannot. Both hold at once, which is the separation of planes of
  Section~\ref{sec:planes}.
\item Moral agency is two theories of the common good over one act. What reads as struggle is their
  simultaneous evaluation of it.
\end{enumerate}

These are the three marks the material of Section~\ref{sec:quantity} has to carry, and the passage shows
that a text can carry all three with none of them named, which is why the analysis needs annotated
material and not a search for words. What the appendix does not show is that a model has these
functions. The text is human and the levels in it are human; how far a Transformer has generalized them
is what Section~\ref{sec:degree} measures, and the profile of Section~\ref{sec:profile} admits a
deficit, a shift or a dropped axis on any component. The passage fixes what to look for and not what
will be found.
